\section{Benchmarking Against Published Observations}
\label{sec:validation}

% =============================================================================
% 4.1  Comparison with published data
% =============================================================================
\subsection{Consistency checks against published resonance data}
\label{sec:published_data}

We benchmark the forward model against published acoustic measurements
on watermelons and other fruits.  These comparisons are
\emph{consistency checks}, not controlled experimental validation:
the published studies report resonant frequencies for fruits of varied
and often incompletely documented geometry and ripeness, and the model
parameters used here are assumed from literature ranges rather than
measured on the same specimens.  We therefore
construct model predictions using the reported fruit dimensions where
available, supplemented by canonical material properties
(\emph{assumed from literature}).

\subsubsection{Yamamoto et al.\ (1980): watermelon tap-tones}

Yamamoto et al.~\cite{Yamamoto1980} measured tap-tone frequencies for
intact watermelons spanning a range of sizes and ripeness states.  They
reported resonant frequencies in the range
\SIrange{80}{180}{\hertz} for fruit with equatorial diameters of
\SIrange{20}{34}{\centi\metre}.  We evaluate the forward model for
three representative sizes at the ripe modulus
($E_{\mathrm{rind}} = \SI{50}{\mega\pascal}$) and three at an
unripe modulus ($E_{\mathrm{rind}} = \SI{150}{\mega\pascal}$),
using geometry consistent with their reported dimensions
(Table~\ref{tab:yamamoto}).

\begin{table}[t]
  \centering
  \caption{Comparison of model predictions with the frequency ranges
    reported by Yamamoto et al.~\cite{Yamamoto1980}.  Model geometry is
    chosen to match reported fruit sizes (\emph{measured from literature});
    material properties are set to
    canonical ripe ($E = \SI{50}{\mega\pascal}$) or unripe
    ($E = \SI{150}{\mega\pascal}$) values (\emph{assumed from literature
    range}).  This is a consistency check, not a point-by-point
    validation.}
  \label{tab:yamamoto}
  \begin{tabular}{lcccccc}
    \toprule
    Description & $a$ [m] & $c$ [m] & $h$ [mm] & $E$ [MPa]
      & $f_2^{\mathrm{model}}$ [Hz] & Yamamoto range [Hz] \\
    \midrule
    Small, ripe     & 0.110 & 0.100 & 13 & 50  & 132 & \multirow{3}{*}{80--120} \\
    Medium, ripe    & 0.140 & 0.120 & 15 & 50  & 102 & \\
    Large, ripe     & 0.170 & 0.140 & 17 & 50  &  83 & \\[3pt]
    Small, unripe   & 0.110 & 0.100 & 15 & 150 & 245 & \multirow{3}{*}{120--180} \\
    Medium, unripe  & 0.140 & 0.120 & 18 & 150 & 192 & \\
    Large, unripe   & 0.170 & 0.140 & 20 & 150 & 155 & \\
    \bottomrule
  \end{tabular}
\end{table}

Three of six predicted frequencies fall within the reported
experimental range (Table~\ref{tab:yamamoto}): the medium and large
ripe cases (\SI{102}{\hertz} and \SI{83}{\hertz}, both within
\SIrange{80}{120}{\hertz}) and the large unripe case
(\SI{155}{\hertz}, within \SIrange{120}{180}{\hertz}).  The remaining
three predictions systematically overshoot the upper experimental
bound---the small ripe prediction (\SI{132}{\hertz}) exceeds the
ripe range by \SI{12}{\hertz}, while the small and medium unripe
predictions (\SI{245}{\hertz} and \SI{192}{\hertz}) exceed the
unripe range by \SI{65}{\hertz} and \SI{12}{\hertz}, respectively.
The mean absolute deviation from the nearest experimental bound is
\SI{15}{\hertz} across all six cases (zero for points within range).
The systematic overprediction for smaller and stiffer fruit suggests
that the canonical moduli assumed here
($E = \SI{150}{\mega\pascal}$ for unripe) may exceed the actual
values for the cultivars studied by Yamamoto et al.\ in 1980, or
that the equivalent-sphere approximation introduces a size-dependent
bias.  Nevertheless, the qualitative scaling---larger fruit yields
lower frequency, stiffer rind yields higher frequency---is captured
correctly, and the model reproduces the correct order of magnitude
across the full size--ripeness matrix.

\subsubsection{Sadrnia et al.\ (2008): vibration-based quality prediction in watermelon}

Sadrnia et al.~\cite{Sadrnia2008Vibration} analysed vibration signals
from intact watermelons as predictors of quality.  Their study is
primarily empirical: it reports spectral features and predictive
performance rather than the full geometric and material information
needed for a point-by-point inversion.  A strict quantitative overlay is
therefore not possible from the published data alone.  Even so, their
central observation---that low-frequency structural vibration of the
intact fruit carries ripeness information---is entirely consistent with
the present model, which identifies the informative observable as the
lowest flexural resonance of a fluid-loaded shell.

\subsubsection{De Belie et al.\ (2000): cross-species consistency check (pear)}

De Belie et al.~\cite{DeBelieetal2000} measured acoustic impulse
response frequencies for Conference pears, not watermelons, and
established an empirical correlation between the resonant frequency and
a stiffness index~$FI$.  We include this comparison deliberately as a
\emph{cross-species consistency check} rather than as a substitute for watermelon
validation.
Although pears are smaller and have a different aspect ratio ($\zeta
\approx 0.6$--$0.7$) than watermelons, they are also well described as
fluid-filled shells.  The existence of a monotonic frequency--stiffness
relationship in pears is qualitatively consistent with the model's
prediction that $f_2 \propto \sqrt{E_{\mathrm{rind}}}$.

\subsubsection{Chen and De Baerdemaeker (1993): cross-species consistency check (pineapple)}

Chen and De Baerdemaeker~\cite{ChenDeBaerdemaeker1993} performed modal
analysis on pineapples, identifying resonant modes that correlated with
flesh firmness.  Their measured frequencies
($\sim\SIrange{100}{300}{\hertz}$) are consistent with the general
prediction of the shell model when pineapple-appropriate geometry and
rind stiffness are used.  This cross-species consistency supports the
suggestion that the equivalent-sphere shell framework may generalise beyond
watermelons, although a rigorous validation for each species would
require purpose-designed experiments.

% =============================================================================
% 4.2  Sensitivity to geometric uncertainty
% =============================================================================
\subsection{Effect of geometric uncertainty}
\label{sec:geom_uncertainty}

In practice, fruit geometry is measured approximately---for example, with
callipers, from a photograph, or by mass--density inference.  We quantify
the effect of geometric measurement error on the inverted modulus through
a Monte Carlo simulation.  The true tap-tone frequency~$f_2$ is computed
from the canonical ripe parameters
($E_{\mathrm{rind}} = \SI{50}{\mega\pascal}$, yielding
$f_2 = \SI{89.9}{\hertz}$), and then the inversion is performed
10\,000~times with each geometric parameter ($a$, $c$, $h$) independently
perturbed by a uniform $\pm\SI{10}{\percent}$ error.

Table~\ref{tab:geom_uncertainty} summarises the results.  The inverted
modulus is unbiased
($\overline{E}_{\mathrm{inv}} = \SI{50.3}{\mega\pascal}$) but exhibits
a standard deviation of \SI{6.8}{\mega\pascal}
(CoV~$= \SI{13.5}{\percent}$), with a \SI{95}{\percent} confidence
interval of $[\SI{38.5}{\mega\pascal}, \SI{64.1}{\mega\pascal}]$.  This
level of uncertainty is acceptable for ripeness classification: the
$\pm\SI{10}{\percent}$ geometric error maps to a $\pm\SI{14}{\percent}$
modulus uncertainty, which is smaller than the factor-of-four modulus
change between ripeness categories (Table~\ref{tab:forward_results}).

\begin{table}[t]
  \centering
  \caption{Effect of $\pm\SI{10}{\percent}$ geometric measurement
    uncertainty on the inverted rind modulus, evaluated by Monte Carlo
    simulation ($n = 10{,}000$) for two ripeness stages.}
  \label{tab:geom_uncertainty}
  \begin{tabular}{lccccc}
    \toprule
    Stage & $E_{\mathrm{true}}$ [MPa] &
      $\overline{E}_{\mathrm{inv}}$ [MPa] &
      $\sigma_E$ [MPa] & CoV [\%] &
      \SI{95}{\percent} CI [MPa] \\
    \midrule
    Ripe   &  50 &  50.3 &  6.8 & 13.5 & [38.5, 64.1] \\
    Unripe & 200 & 201.3 & 27.7 & 13.7 & --- \\
    \bottomrule
  \end{tabular}
\end{table}

We also examine the individual sensitivity of $f_2$ to each geometric
parameter.  For the canonical ripe watermelon
($f_2 = \SI{89.9}{\hertz}$), a $\pm\SI{10}{\percent}$ perturbation in
each parameter yields:
\begin{itemize}
  \item $a$: $\Delta f_2 \in [\SI{-8.8}{\percent}, \SI{+10.7}{\percent}]$
    --- the strongest geometric sensitivity, consistent with
    $f_2 \propto a^{-3/2}$;
  \item $c$: $\Delta f_2 \in [\SI{-4.5}{\percent}, \SI{+5.2}{\percent}]$;
  \item $h$: $\Delta f_2 \in [\SI{-4.6}{\percent}, \SI{+4.4}{\percent}]$;
  \item $\rho_f$: $\Delta f_2 \in [\SI{-3.8}{\percent}, \SI{+4.3}{\percent}]$;
  \item $\rho_w$: $\Delta f_2 \in [\SI{-0.9}{\percent}, \SI{+0.9}{\percent}]$ --- negligible.
\end{itemize}
The semi-major axis~$a$ dominates, consistent with the Sobol analysis
(\S\ref{sec:sobol}).  In a practical measurement scenario, even a crude
estimate of the equatorial diameter (obtainable from a single photograph)
reduces the dominant source of geometric uncertainty.

\begin{figure}[t]
  \centering
  \includegraphics[width=0.7\textwidth]{fig_inversion_noise}
  \caption{Effect of frequency measurement noise on inverted rind modulus,
    assessed by 1\,000 Monte Carlo trials at each noise level.  The
    condition number of the inversion with respect to~$f$ is
    approximately~2 (since $E \propto f_2^2$, the condition number with
    respect to~$f^2$ is unity): \SI{2}{\percent} frequency noise yields
    $\approx\SI{4}{\percent}$ modulus error.  Violin plots show the full
    distribution; box plots indicate interquartile range.}
  \label{fig:inversion_noise}
\end{figure}
